% Editorial sheet fixing one ordered-superposition presentation. Completeness
% depends on the ordering, eligibility, redundancy, and fair saturation clauses.
\calculussheet{sheet:ordered-superposition}{Ordered Superposition with Selection}{Language: first-order clauses whose literals are equations $s\approx t$ or disequations $s\not\approx t$; $\succ$ is a complete simplification ordering on ground terms extended to literals/clauses; a selection function selects negative literals, and otherwise eligible literals are maximal; clauses are considered modulo variable renaming; equality has identity semantics.}
\noindent\emph{Controlling source:} \cite{BachmairGanzinger2001,NieuwenhuisRubio2001}.
\hbrule{rule:superposition}{Superposition}{C\lor l\approx r\qquad D\lor s[l']\mathrel{\diamond}t}{(C\lor D\lor s[r]\mathrel{\diamond}t)\sigma}{$\diamond\in\{\approx,\not\approx\}$; $l'$ is nonvariable; $\sigma$ is an MGU of $l,l'$; $l\sigma\succ r\sigma$ and $s[l']\sigma\succ t\sigma$; the source equality and target literal are eligible under the fixed selection/ordering policy}
\hbrule{rule:equality-resolution}{Equality Resolution}{C\lor s\not\approx t}{C\sigma}{$\sigma$ is an MGU of $s,t$ and the disequality is eligible}
\hbrule{rule:equality-factoring}{Equality Factoring}{C\lor s\approx t\lor u\approx v}{(C\lor t\not\approx v\lor u\approx v)\sigma}{$\sigma$ is an MGU of $s,u$; $(s\approx t)\sigma$ is maximal in the instantiated premise and no literal is selected; $s\sigma\not\preceq t\sigma$ and $u\sigma\not\preceq v\sigma$}
